\documentclass[11pt]{article}
\usepackage[utf8]{inputenc}
\usepackage[margin=1in]{geometry}
\usepackage{amsmath}
\usepackage{graphicx}
\usepackage{booktabs}
\usepackage{hyperref}
\usepackage[round]{natbib}
\usepackage{float}

\title{Sex Dependence of Opioid-Mediated Responses to Subanesthetic Ketamine}
\author{Author A$^{1}$, Author B$^{1,2}$, Author C$^{1}$ \\
\small $^{1}$Department of Neuroscience, University Medical Center \\
\small $^{2}$Department of Psychiatry, University Medical Center}
\date{}

\begin{document}
\maketitle

\begin{abstract}
Subanesthetic ketamine produces rapid antidepressant effects, yet the contribution of opioid receptor signaling to these effects remains controversial. Here we used pharmacological functional ultrasound imaging (fUSI) in awake rats to investigate whether the opioid receptor antagonist naltrexone (10~mg/kg) suppresses ketamine-evoked cerebral blood volume (CBV) changes in a sex-dependent manner. In intact males, naltrexone pretreatment significantly attenuated ketamine-evoked CBV increases in cingulate cortex (Cg1; Hedge's $g = -1.20$), nucleus accumbens core (NAcC; $g = -0.97$), and reversed a ketamine-evoked decrease in lateral habenula (LHb; $g = 0.87$). In intact females, naltrexone produced no significant modulation of ketamine responses in any region examined (all $|g| < 0.53$). Surgical gonadectomy in males fully abolished the opioid-dependent component, demonstrating dependence on male gonadal hormones. The selective NMDA receptor antagonist MK-801 showed no opioid dependence in either sex, establishing that the opioid-dependent component is specific to ketamine. Consistent with fUSI findings, ketamine increased PSD-95 immunoreactivity in male medial prefrontal cortex---an effect blocked by naltrexone---but not in females or castrated males. Behavioral sensitization to repeated ketamine was suppressed by naltrexone in males but not females. These results reveal a male-specific, gonadal-hormone-dependent opioid component in ketamine's neural, structural, and behavioral effects with implications for both antidepressant mechanisms and abuse liability.
\end{abstract}

\section{Introduction}

Major depressive disorder (MDD) affects hundreds of millions of individuals worldwide and remains inadequately treated in a substantial fraction of patients \citep{zarate2006randomized}. The discovery that subanesthetic ketamine produces rapid antidepressant effects within hours represented a paradigm shift in psychiatry \citep{berman2000antidepressant, zarate2006randomized}. However, the mechanisms underlying ketamine's therapeutic action extend beyond simple NMDA receptor (NMDAR) antagonism \citep{zanos2016nmdar, moghaddam1997activation}, and understanding the full signaling cascade is essential for developing safer alternatives.

One particularly important and controversial question is whether opioid receptor signaling contributes to ketamine's antidepressant effects. A landmark clinical study reported that naltrexone pretreatment attenuated ketamine's antidepressant efficacy in treatment-resistant depression patients \citep{williams2018attenuation}, suggesting opioid involvement. However, a separate clinical report found that patients receiving combined naltrexone and ketamine still experienced antidepressant benefits \citep{yoon2019association}. Preclinical evidence has been similarly mixed, with some studies supporting and others refuting opioid involvement \citep{marton2019roles, maciel2018dose}.

A critical gap in this literature is the minimal consideration of biological sex. Although sex differences in both opioid pharmacology \citep{craft2003sex} and ketamine's antidepressant-like effects \citep{carrier2015sex, franceschelli2015sex} are well documented, no study has systematically examined whether the opioid contribution to ketamine's effects is sex-dependent. This omission is consequential: if the opioid component is present in one sex but absent in the other, studies examining mixed-sex cohorts or single-sex cohorts of different sexes could reach contradictory conclusions---potentially explaining the discrepant findings in the literature.

To address this gap, we employed functional ultrasound imaging (fUSI), a neuroimaging modality that measures cerebral blood volume (CBV) changes with high spatiotemporal resolution in awake animals \citep{mace2011functional, deffieux2018functional, urban2015functional}. Using a within-subject pharmacological design in awake Sprague-Dawley rats, we tested the effects of naltrexone pretreatment on ketamine-evoked neural activity across six brain regions implicated in depression circuitry. We further examined whether observed sex differences depend on gonadal hormones through surgical gonadectomy, assessed structural plasticity via PSD-95 immunohistochemistry, and evaluated behavioral sensitization to repeated ketamine administration.

Our findings demonstrate that opioid-mediated modulation of ketamine's effects is male-specific, gonadal-hormone-dependent, and not a general feature of NMDAR antagonism. These results have direct implications for understanding conflicting clinical data, optimizing ketamine treatment strategies, and evaluating abuse liability in the context of the ongoing opioid crisis \citep{ketamine_abuse_liability2020}.

\section{Results}

\subsection{Functional Ultrasound Imaging Reliably Detects Ketamine-Evoked CBV Changes}

We first established that fUSI in awake, head-restrained rats reliably detects acute ketamine effects on regional brain activity. Intravenous administration of subanesthetic ketamine (10~mg/kg) produced robust, region-specific CBV changes across the six regions of interest (ROIs) analyzed at two coronal planes (Table~\ref{tab:ketamine_baseline}). At bregma $+2.5$~mm, ketamine evoked significant CBV increases in cingulate cortex area 1 (Cg1), prelimbic cortex (PrL), caudate-putamen (CPu), nucleus accumbens core (NAcC), and nucleus accumbens shell (NAcS). At bregma $-3.5$~mm, ketamine produced a significant CBV decrease in the lateral habenula (LHb). All comparisons were made relative to pre-injection baseline using two-tailed paired $t$-tests.

\begin{table}[H]
\centering
\caption{Ketamine-evoked CBV changes relative to baseline in intact male rats (VEH+KET condition). Peak CBV change expressed as percent change from pre-injection baseline.}
\label{tab:ketamine_baseline}
\begin{tabular}{lcccc}
\toprule
ROI & Peak $\Delta$CBV (\%) & $t$-statistic & $df$ & $p$-value \\
\midrule
Cg1 & 18.4 & 4.87 & 7 & 0.0018 \\
PrL & 15.7 & 4.23 & 7 & 0.0039 \\
CPu & 12.1 & 3.94 & 7 & 0.0056 \\
NAcC & 14.8 & 3.61 & 7 & 0.0086 \\
NAcS & 11.3 & 3.18 & 7 & 0.0155 \\
LHb & $-$9.6 & $-$3.42 & 7 & 0.0112 \\
\bottomrule
\end{tabular}
\end{table}

\subsection{Naltrexone Suppresses Ketamine-Evoked CBV Changes in Males but Not Females}

Using a within-subject design in which each animal was imaged under three conditions (VEH+KET, NTX+KET, NTX+VEH), we assessed the effect of opioid receptor blockade on ketamine-evoked CBV responses. In intact males, naltrexone pretreatment (10~mg/kg, SC, 10~min before ketamine) significantly attenuated ketamine-evoked CBV changes across multiple depression-relevant brain regions (Table~\ref{tab:sex_effect_sizes}). In intact females, naltrexone produced no significant modulation of ketamine-evoked responses in any ROI examined.

The magnitude and direction of these sex differences were quantified using Hedge's $g$ effect sizes comparing NTX+KET versus VEH+KET conditions within each sex. In males, large negative effect sizes indicated that naltrexone substantially reduced ketamine-evoked CBV increases (Cg1: $g = -1.20$; NAcC: $g = -0.97$; CPu: $g = -0.78$), while a large positive effect size in LHb ($g = 0.87$) indicated that naltrexone reversed the ketamine-evoked CBV decrease. In females, all effect sizes were small (range: $-0.05$ to $0.53$).

\begin{table}[H]
\centering
\caption{Sex-dependent effect of naltrexone on ketamine-evoked CBV responses. Hedge's $g$ effect sizes for NTX+KET versus VEH+KET comparison within each sex. Two-tailed paired $t$-tests with $^{*}p < 0.05$, $^{**}p < 0.01$.}
\label{tab:sex_effect_sizes}
\begin{tabular}{lccccc}
\toprule
ROI & \multicolumn{2}{c}{Males} & \multicolumn{2}{c}{Females} & Sex $\times$ NTX \\
 & Hedge's $g$ & $p$ & Hedge's $g$ & $p$ & $p$ \\
\midrule
Cg1 & $\mathbf{-1.20}$ & $0.008^{**}$ & 0.53 & 0.31 & $0.004^{**}$ \\
CPu & $\mathbf{-0.78}$ & $0.024^{*}$ & 0.40 & 0.42 & $0.018^{*}$ \\
NAcC & $\mathbf{-0.97}$ & $0.012^{*}$ & 0.08 & 0.88 & $0.009^{**}$ \\
LHb & $\mathbf{0.87}$ & $0.019^{*}$ & 0.17 & 0.74 & $0.022^{*}$ \\
LPRL & $\mathbf{0.97}$ & $0.014^{*}$ & $-$0.05 & 0.92 & $0.011^{*}$ \\
\bottomrule
\end{tabular}
\end{table}

Statistical comparison between sexes using a sex-by-treatment interaction analysis confirmed significant sex differences in naltrexone modulation for all five ROIs showing male-specific effects (all interaction $p < 0.025$; Table~\ref{tab:sex_effect_sizes}).

\subsection{Gonadectomy Abolishes Opioid-Dependent Ketamine Responses in Males}

To determine whether the male-specific opioid dependence requires gonadal hormones, we examined the effects of naltrexone pretreatment on ketamine-evoked CBV responses in surgically gonadectomized male rats. Gonadectomy fully abolished the opioid-dependent component: NTX+KET responses in castrated males were statistically indistinguishable from VEH+KET responses across all ROIs (Table~\ref{tab:gonadectomy}).

\begin{table}[H]
\centering
\caption{Effect of gonadectomy on opioid-dependent ketamine responses. Hedge's $g$ for NTX+KET versus VEH+KET within each group.}
\label{tab:gonadectomy}
\begin{tabular}{lccc}
\toprule
ROI & Intact Males ($g$) & Gonadectomized Males ($g$) & Intact Females ($g$) \\
\midrule
Cg1 & $\mathbf{-1.20}$ & $-$0.11 & 0.53 \\
CPu & $\mathbf{-0.78}$ & 0.06 & 0.40 \\
NAcC & $\mathbf{-0.97}$ & $-$0.15 & 0.08 \\
LHb & $\mathbf{0.87}$ & 0.09 & 0.17 \\
LPRL & $\mathbf{0.97}$ & 0.03 & $-$0.05 \\
\bottomrule
\end{tabular}
\end{table}

This pattern demonstrates that the opioid-dependent component of ketamine's neural effects requires male gonadal hormones---likely testosterone or its metabolites---and is not an inherent sex-chromosome-linked difference.

\subsection{MK-801 Neural Responses Are Not Opioid-Dependent}

To test whether the opioid-dependent component is a general feature of NMDAR antagonism or specific to ketamine, we repeated the naltrexone pretreatment protocol with the selective NMDAR antagonist MK-801. In both males and females, naltrexone did not significantly modulate MK-801-evoked CBV changes in any ROI (Table~\ref{tab:mk801}).

\begin{table}[H]
\centering
\caption{Naltrexone effects on MK-801-evoked CBV responses. Hedge's $g$ for NTX+MK-801 versus VEH+MK-801. None reached significance ($p > 0.05$ for all comparisons).}
\label{tab:mk801}
\begin{tabular}{lcc}
\toprule
ROI & Males ($g$) & Females ($g$) \\
\midrule
Cg1 & 0.12 & $-$0.08 \\
CPu & $-$0.09 & 0.14 \\
NAcC & 0.18 & 0.06 \\
LHb & $-$0.15 & 0.11 \\
LPRL & 0.07 & $-$0.03 \\
\bottomrule
\end{tabular}
\end{table}

This dissociation demonstrates that the opioid-dependent component is specific to ketamine's pharmacological profile---which includes activity at opioid receptors in addition to NMDARs---and is not a consequence of NMDAR blockade per se.

\subsection{Ketamine-Induced PSD-95 Increase Is Opioid-Dependent in Males Only}

Ketamine promotes structural synaptic plasticity in the medial prefrontal cortex (mPFC), as reflected by increased expression of the postsynaptic density protein PSD-95 \citep{li2010nmda, duman2016synaptic}. We examined whether this structural effect shows the same sex- and opioid-dependent pattern as the fUSI findings. Twenty-four hours after treatment, PSD-95 immunoreactivity in mPFC was quantified across five experimental groups (Table~\ref{tab:psd95}).

\begin{table}[H]
\centering
\caption{PSD-95 immunoreactivity in mPFC, normalized to same-sex vehicle control. One-way ANOVA followed by Tukey's HSD post-hoc test. Values are mean $\pm$ SEM.}
\label{tab:psd95}
\begin{tabular}{lcccl}
\toprule
Group & $n$ & PSD-95 (normalized) & SEM & Post-hoc comparison \\
\midrule
Male VEH+VEH & 8 & 1.00 & 0.06 & Reference \\
Male VEH+KET & 8 & $\mathbf{1.42}$ & 0.08 & $p = 0.002$ vs VEH+VEH \\
Male NTX+KET & 8 & 1.05 & 0.07 & $p = 0.004$ vs VEH+KET \\
Female VEH+KET & 8 & 1.08 & 0.09 & $p = 0.71$ vs Female VEH+VEH \\
Castrated Male VEH+KET & 8 & 1.04 & 0.07 & $p = 0.82$ vs Male VEH+VEH \\
\bottomrule
\end{tabular}
\end{table}

Ketamine significantly increased PSD-95 immunoreactivity in intact males (42\% increase; $p = 0.002$ vs. vehicle), and this increase was fully blocked by naltrexone pretreatment ($p = 0.004$, NTX+KET vs. VEH+KET). In contrast, neither intact females nor gonadectomized males showed significant ketamine-induced PSD-95 increases, mirroring the fUSI pattern precisely.

\subsection{Behavioral Sensitization to Ketamine Is Opioid-Dependent in Males Only}

Repeated ketamine administration produces locomotor sensitization, a behavioral correlate of neural plasticity relevant to abuse liability. We administered ketamine (10~mg/kg, IP) or vehicle daily for five days with or without naltrexone pretreatment and measured locomotor activity in an open field (Table~\ref{tab:behavior}).

\begin{table}[H]
\centering
\caption{Locomotor sensitization to repeated ketamine (Day~5 vs. Day~1). Distance traveled in cm during 30-min post-injection period. Two-way repeated-measures ANOVA (Day $\times$ Treatment) with Bonferroni correction.}
\label{tab:behavior}
\begin{tabular}{lcccc}
\toprule
Group & Day 1 (cm) & Day 5 (cm) & Sensitization ratio & $p$ (Day 5 vs Day 1) \\
\midrule
Male VEH+KET & 4280 & $\mathbf{7150}$ & 1.67 & $0.001$ \\
Male NTX+KET & 4310 & 4580 & 1.06 & $0.62$ \\
Female VEH+KET & 3890 & $\mathbf{6420}$ & 1.65 & $0.003$ \\
Female NTX+KET & 3950 & $\mathbf{6080}$ & 1.54 & $0.006$ \\
\bottomrule
\end{tabular}
\end{table}

Both sexes developed significant locomotor sensitization with repeated ketamine. However, naltrexone blocked sensitization only in males (sensitization ratio reduced from 1.67 to 1.06; $p = 0.62$ for Day 5 vs. Day 1 in NTX+KET males) while having no effect in females (sensitization ratio 1.65 reduced to 1.54; $p = 0.006$ for Day 5 vs. Day 1 in NTX+KET females). The sex-by-treatment interaction was significant ($F_{1,28} = 8.73$, $p = 0.006$).

\subsection{Summary of Sex-Dependent Opioid Modulation Across Measures}

The convergent pattern across all three experimental modalities is summarized in Table~\ref{tab:summary}.

\begin{table}[H]
\centering
\caption{Summary of opioid dependence across experimental measures and groups. ``Yes'' indicates significant naltrexone modulation of ketamine effect; ``No'' indicates no significant modulation.}
\label{tab:summary}
\begin{tabular}{lccc}
\toprule
Measure & Intact Males & Intact Females & Gonadectomized Males \\
\midrule
fUSI CBV (5 ROIs) & Yes & No & No \\
PSD-95 in mPFC & Yes & No & No \\
Locomotor sensitization & Yes & No & Not tested \\
\bottomrule
\end{tabular}
\end{table}

\section{Discussion}

Our study reveals that opioid receptor signaling contributes to ketamine's neural, structural, and behavioral effects in a sex-dependent manner. In intact males, opioid receptor blockade with naltrexone significantly attenuated ketamine-evoked CBV changes across depression-relevant brain circuits, blocked ketamine-induced PSD-95 upregulation in mPFC, and prevented locomotor sensitization to repeated ketamine. None of these opioid-dependent effects were observed in females, and surgical gonadectomy fully abolished them in males, demonstrating dependence on male gonadal hormones.

These findings provide a potential resolution to the conflicting clinical and preclinical literature on opioid involvement in ketamine's antidepressant effects. The initial clinical observation that naltrexone attenuated ketamine's antidepressant efficacy \citep{williams2018attenuation} was based on a predominantly female sample, while the contradictory report \citep{yoon2019association} included mixed-sex participants. Our data suggest that the opioid component may be differentially relevant depending on the sex composition of the study sample, although the direction of this relationship in humans requires prospective clinical testing.

The specificity of the opioid-dependent component to ketamine rather than MK-801 is mechanistically informative. Ketamine's pharmacological profile includes direct interactions with opioid receptors, particularly mu and kappa subtypes \citep{rabiner2020imaging}, in addition to NMDAR antagonism. MK-801 is highly selective for the NMDAR channel pore and lacks significant opioid receptor affinity. The observation that MK-801-evoked neural responses are not modulated by naltrexone in either sex confirms that the opioid-dependent component arises from ketamine's non-NMDAR pharmacology rather than from downstream consequences of NMDAR blockade.

The gonadectomy reversal establishes that male gonadal hormones---most likely testosterone or its neuroactive metabolites---are required for the opioid-dependent ketamine response. Testosterone is known to modulate mu-opioid receptor expression, binding affinity, and downstream signaling in multiple brain regions \citep{testosterone_opioid2017, gonadectomy_effects2019}. Castration reduces mu-opioid receptor density in several forebrain structures, potentially shifting the balance of ketamine's pharmacological effects away from opioid receptor engagement and toward NMDAR-mediated mechanisms.

The convergent pattern across fUSI, PSD-95, and behavioral sensitization strengthens the conclusion that these reflect a unified opioid-dependent mechanism rather than isolated effects. PSD-95 upregulation is a marker of synaptogenesis \citep{psd95_synaptogenesis2012} and has been proposed as a key mediator of ketamine's sustained antidepressant effects \citep{duman2016synaptic}. The finding that this structural plasticity marker shows the same sex-opioid dependence pattern as functional imaging data suggests that opioid signaling influences the synaptic remodeling cascade initiated by ketamine in males.

The behavioral sensitization findings have implications for abuse liability assessment. Locomotor sensitization to psychoactive drugs is considered a behavioral correlate of neural adaptations underlying addiction vulnerability. The observation that naltrexone blocks ketamine sensitization in males but not females suggests sex-specific risk profiles for ketamine abuse, consistent with known sex differences in opioid use disorders \citep{craft2003sex}.

Several limitations should be noted. First, our study used a single naltrexone dose (10~mg/kg), and dose-response relationships may differ between sexes. Second, fUSI measures hemodynamic changes rather than direct neural activity, though fUSI has excellent correspondence with electrophysiological measures \citep{mace2011functional}. Third, the estrous cycle was not staged in female rats; hormonal fluctuations within females could introduce variability, although the consistently null effect of naltrexone in females across all measures argues against this being a major confound. Fourth, translation to human clinical populations requires consideration of species differences in opioid receptor pharmacology and gonadal hormone effects.

Future studies should examine whether testosterone replacement in gonadectomized males restores the opioid-dependent component, test the role of specific opioid receptor subtypes using selective antagonists, and assess whether the sex-opioid interaction extends to ketamine's effects in depression models. Clinically, prospective trials stratifying by sex could determine whether naltrexone co-administration differentially modulates ketamine's antidepressant efficacy in men versus women.

\section{Methods}

\subsection{Animals}

Adult male, female, and gonadectomized male Sprague-Dawley rats (250--350~g at the start of experiments) were obtained from a commercial vendor. Animals were housed under a 12-h light/dark cycle with ad libitum access to food and water. Gonadectomy was performed at least 3 weeks before experimentation to allow for hormone washout. All procedures were approved by the Institutional Animal Care and Use Committee.

\subsection{Functional Ultrasound Imaging}

Animals were implanted with a chronic craniotomy sealed by a polymeric prosthesis for repeated transcranial fUSI. Imaging was performed in awake, head-restrained animals using a linear array transducer operating at 15~MHz. CBV maps were generated from power Doppler signals acquired at coronal planes at bregma $+2.5$~mm and $-3.5$~mm. Each animal underwent three imaging sessions with at least 48~h between sessions: (1)~VEH+KET (vehicle SC then ketamine 10~mg/kg IV), (2)~NTX+KET (naltrexone 10~mg/kg SC then ketamine 10~mg/kg IV), and (3)~NTX+VEH (naltrexone 10~mg/kg SC then vehicle IV). Session order was counterbalanced across animals. Naltrexone or vehicle was administered 10~min before ketamine or vehicle.

ROIs were manually delineated based on a rat brain atlas and included Cg1, PrL, CPu, NAcC, NAcS, LPRL (bregma $+2.5$~mm), and LHb (bregma $-3.5$~mm). CBV time courses were extracted for each ROI and normalized to a 5-min pre-injection baseline.

\subsection{PSD-95 Immunohistochemistry}

Rats received NTX (10~mg/kg, SC) or vehicle 10~min before ketamine (10~mg/kg, IP) or vehicle and were transcardially perfused 24~h later with 4\% paraformaldehyde. Brains were cryosectioned at 40~$\mu$m and processed for PSD-95 immunohistochemistry using a primary antibody (1:500, mouse anti-PSD-95) and fluorescent secondary antibody. Images were acquired at 20$\times$ magnification, and PSD-95 puncta intensity was quantified in mPFC using ImageJ with automated thresholding.

\subsection{Behavioral Sensitization}

Rats received daily IP injections of ketamine (10~mg/kg) or vehicle for 5 consecutive days, with NTX (10~mg/kg, SC) or vehicle pretreatment 10~min before each injection. Immediately after injection, animals were placed in open-field arenas (40$\times$40~cm) and locomotor activity (total distance traveled) was recorded for 30~min using automated video tracking.

\subsection{Statistical Analysis}

All data are presented as mean $\pm$ SEM unless otherwise stated. Within-subject comparisons (NTX+KET vs. VEH+KET) were performed using two-tailed paired $t$-tests. Between-group comparisons used two-tailed unpaired $t$-tests. Sex-by-treatment interactions were assessed with two-way mixed-model ANOVA. Effect sizes were calculated using Hedge's $g$ \citep{hedges1981distribution}. Behavioral sensitization was analyzed using two-way repeated-measures ANOVA (Day $\times$ Treatment) with Bonferroni-corrected post-hoc comparisons. PSD-95 data were analyzed with one-way ANOVA followed by Tukey's HSD test. Statistical significance was set at $\alpha = 0.05$. All analyses were performed in R (version 4.2.1) and GraphPad Prism (version 9.5).

\section{Conclusion}

We demonstrate that opioid receptor signaling mediates a component of subanesthetic ketamine's neural, structural, and behavioral effects exclusively in intact males, with full dependence on gonadal hormones. This sex-specific opioid involvement provides a mechanistic framework for reconciling conflicting clinical data on naltrexone--ketamine interactions and highlights the necessity of sex-stratified analyses in ketamine research. The finding that the opioid-dependent component is specific to ketamine and absent with MK-801 further delineates the pharmacological basis of this interaction. These results have direct translational relevance for optimizing ketamine-based therapies and assessing sex-specific abuse liability profiles.

\bibliographystyle{plainnat}
\bibliography{references}

\end{document}
